\exercisesection{}

\begin{enumerate}
\def\labelenumi{\arabic{enumi})}
\item
  设 E 和 F 是希尔伯特空间，u 是从 E 到 F 的一个连续线性映射。证明 u 是弗雷德霍姆映射，当且仅当 u 的像是闭的且 u 和 \(u^{*}\) 的核是有限维的。
\item
  设 \(p\) 是满足 \(1 \leqslant p < +\infty\) 的实数。设 \(u\) 是 \(\ell_{\mathrm{K}}^{p}(\mathbf{N})\) 的一个对角自同态。证明 \(u\) 是弗雷德霍姆自同态，当且仅当它是里斯自同态。特别地，\(\ell_{\mathrm{K}}^{p}(\mathbf{N})\) 的每个对角弗雷德霍姆自同态的指标为零。
\item
  \begin{enumerate}
  \def\labelenumii{\alph{enumii})}
  \tightlist
  \item
    存在一个巴拿赫空间 E、E 的一个连续自同构 u 以及 E 的一个里斯自同态 v，使得 \(u \circ v\) 不是里斯自同态。甚至存在这样的例子，其中 \(u \circ v - v \circ u\) 是有限秩的。
  \end{enumerate}
\end{enumerate}

\begin{enumerate}
\def\labelenumi{\alph{enumi})}
\setcounter{enumi}{1}
\tightlist
\item
  存在一个巴拿赫空间 E、E 的一个自同构 u 以及一个有限秩自同态 h，使得 \(u + h\) 不是里斯自同态。
\end{enumerate}

\begin{enumerate}
\def\labelenumi{\arabic{enumi})}
\setcounter{enumi}{3}
\tightlist
\item
  设 E 是无限维复希尔伯特空间。
\end{enumerate}

\begin{enumerate}
\def\labelenumi{\alph{enumi})}
\item
  E 上的酉算子群 \(\mathbf{U}(\mathrm{E})\)，赋以由 \(\mathcal{L}(\mathrm{E})\) 的范数诱导的拓扑，是道路连通的。（证明若 u 是酉的，则存在一个埃尔米特算子 v，使得 \(u = \exp(iv)\) 。）
\item
  E 的线性同构群是道路连通的。（利用极分解，参见 I, p. 139, 定义 3。）
\item
  对每个 \(n \in Z\)，E 的指标为 n 的弗雷德霍姆自同态组成的集合是开且连通的。它是非空的。
\end{enumerate}

\begin{enumerate}
\def\labelenumi{\arabic{enumi})}
\setcounter{enumi}{4}
\tightlist
\item
  设 E 和 F 是无限维巴拿赫空间。连续线性映射 \(u: \mathrm{E} \to \mathrm{F}\) 称为严格奇异的，如果不存在无限维的闭子空间 \(\mathrm{F}_1 \subset \mathrm{E}\)，使得 \(u\) 通过过渡到子空间诱导一个从 \(\mathrm{F}_1\) 到 \(u(\mathrm{F}_1)\) 上的同构。
\end{enumerate}

\begin{enumerate}
\def\labelenumi{\alph{enumi})}
\tightlist
\item
  设 u 是从 E 到 F 的一个连续线性映射。下列条件等价：
\end{enumerate}

\begin{enumerate}
\def\labelenumi{(\roman{enumi})}
\item
  映射 \(u\) 是严格奇异的；
\item
  对 E 的每个无限维闭子空间 \(F_{1}\)，不存在实数 \(\delta > 0\)，使得 \(\|u(x)\| \geqslant \delta\|x\|\) 对所有 \(x \in F_{1}\) 成立。
\end{enumerate}

\begin{enumerate}
\def\labelenumi{\alph{enumi})}
\setcounter{enumi}{1}
\item
  设 \(u \colon \mathrm{E} \to \mathrm{F}\) 是一个连续线性映射，使得 \(u\) 的像是闭的且 \(u\) 的核是有限维的。若 \(v \colon \mathrm{E} \to \mathrm{F}\) 是严格奇异的，则 \(u + v\) 的像是闭的且其核是有限维的。
\item
  E 的严格奇异自同态组成的集合是 \(\mathcal{L}(\mathrm{E})\) 的一个闭双边理想。
\item
  每个紧线性映射都是严格奇异的。若 E 和 F 是希尔伯特空间，则每个严格奇异映射 \(E \rightarrow F\) 都是紧的。
\item
  存在巴拿赫空间 E 和 F 以及一个严格奇异但非紧的映射 \(E \rightarrow F\) 。
\end{enumerate}

\(f)\) 设 \(p\) 和 \(r\) 是满足 \(1 \leqslant p < r\) 的实数。每个连续线性映射 \(\ell_{\mathrm{K}}^{r}(\mathbf{N}) \to \ell_{\mathrm{K}}^{p}(\mathbf{N})\) 都是严格奇异的。（利用 III, p. 106 习题 10。）

\begin{enumerate}
\def\labelenumi{\arabic{enumi})}
\setcounter{enumi}{5}
\item
  设 E 是复希尔伯特空间。设 u 是 E 的一个弗雷德霍姆自同态。若 u 是正规的，则 \(\text{ind}(u) = 0\) 。
\item
  设 E 和 F 是巴拿赫空间，u 是 \(\mathcal{L}(\mathrm{E};\mathrm{F})\) 的一个元素。证明 u 是弗雷德霍姆映射，当且仅当下列条件被满足：
\end{enumerate}

\begin{enumerate}
\def\labelenumi{(\roman{enumi})}
\tightlist
\item
  存在 E 上的一个连续半范数 p 和 \(C \geqslant 0\)，使得从 E 到 \((E, p)\) 的分离完备化的典范映射是紧的，且
\end{enumerate}

\[
\| x \| \leqslant \mathrm{C} \| u (x) \| + p (x)
\]

对所有 \(x \in E\) 成立；

\begin{enumerate}
\def\labelenumi{(\roman{enumi})}
\setcounter{enumi}{1}
\tightlist
\item
  存在一个连续半范数 \(q\)（\(\mathrm{F}'\) 上的）和 \(\mathrm{C}' \geqslant 0\)，使得从 \(\mathrm{F}'\) 到 \((\mathrm{F}', q)\) 的分离完备化的典范映射是紧的，且
\end{enumerate}

\[
\| y ^ {\prime} \| \leqslant \mathrm{C} ^ {\prime} \| ^ {t} u (y ^ {\prime}) \| + q (y ^ {\prime})
\]

对所有 \(y' \in F'\) 成立。

\begin{enumerate}
\def\labelenumi{\arabic{enumi})}
\setcounter{enumi}{7}
\tightlist
\item
  设 U 是 C 的一个连通开子集。我们用 \(\mathcal{H}(\mathrm{U})\) 表示 U 上的赋有紧收敛拓扑的全纯函数空间（EVT III, p. 10, 例 \(d\)）。
\end{enumerate}

\begin{enumerate}
\def\labelenumi{\alph{enumi})}
\item
  设 \(a \in \mathcal{H}(U)\) 。乘以 a 是 \(\mathcal{H}(U)\) 的一个弗雷德霍姆自同态，当且仅当 a 的零点集是有限的。然后计算它的指标。
\item
  设 \(\mathcal{U} = (\mathrm{U}_{i})_{i \in \mathrm{I}}\) 是 U 的一个由非空连通开集组成的覆盖。我们令
\end{enumerate}

\[
\begin{array}{l} \mathrm{I} _ {2} = \{(i, j) \in \mathrm{I} \times \mathrm{I} | \mathrm{U} _ {i} \cap \mathrm{U} _ {j} \neq \varnothing \}, \\ \mathrm{I} _ {3} = \{(i, j, k) \in \mathrm{I} \times \mathrm{I} \times \mathrm{I} | \mathrm{U} _ {i} \cap \mathrm{U} _ {j} \cap \mathrm{U} _ {k} \neq \varnothing \}. \end{array}
\]

我们考虑向量空间

\[
\begin{array}{l} \mathrm{Z} ^ {1} (\mathcal {U}; \mathbf {C}) = \{(c _ {i, j}) \in \mathbf {C} ^ {\mathrm{I} _ {2}} | c _ {i j} + c _ {j k} + c _ {k i} = 0 \text { for all } (i, j, k) \in \mathrm{I} _ {3} \} \\ \mathrm{B} ^ {1} (\mathcal {U}; \mathbf {C}) = \{(c _ {i, j}) \in \mathbf {C} ^ {\mathrm{I} _ {2}} | \text { there exists } c ^ {\prime} \in \mathbf {C} ^ {\mathrm{I}} \text { such that } \end{array}
\]

\[
c _ {i j} = c _ {i} ^ {\prime} - c _ {j} ^ {\prime} \text {   for all   } (i, j) \in \mathrm{I} _ {2} \}.
\]

我们有 \(\mathrm{B}^1 (\mathcal{U};\mathbf{C})\subset \mathrm{Z}^1 (\mathcal{U};\mathbf{C})\) 。我们令

\[
\mathrm{H} ^ {1} (\mathcal {U}; \mathbf {C}) = \mathrm{Z} ^ {1} (\mathcal {U}; \mathbf {C}) / \mathrm{B} ^ {1} (\mathcal {U}; \mathbf {C}).
\]

设 \(\mathcal{V}=(\mathrm{V}_{\alpha})_{\alpha\in\mathrm{A}}\) 是 U 的一个由非空连通开集组成且比 U 更细的覆盖，并设 \(\varrho\colon A\to I\) 是一个映射，使得 \(V_{\alpha}\subset U_{\varrho(\alpha)}\) 对所有 \(\alpha\in A\) 成立。

映射 \((c_{ij})_{i,j} \mapsto (c_{\varrho(\alpha)\varrho(\beta)})_{\alpha,\beta}\) 诱导一个不依赖于 \(\varrho\) 的选取的线性映射，从 \(\mathrm{H}^1(\mathcal{U};\mathbf{C})\) 到 \(\mathrm{H}^1(\mathcal{V};\mathbf{C})\) 。

我们用 \(H^{1}(U; \mathbf{C})\) 表示当 U 取遍 U 的由非空连通开集组成的覆盖时所得到的赋有这些线性映射的向量空间 \(H^{1}(\mathcal{U}; \mathbf{C})\) 的归纳极限。

\begin{enumerate}
\def\labelenumi{\alph{enumi})}
\setcounter{enumi}{2}
\tightlist
\item
  设 D 是 \(\mathcal{H}(\mathrm{U})\) 的由 \(Df = f'\) 定义的自同态。证明存在一个从 D 的余核到 \(\mathrm{H}^1 (\mathrm{U};\mathbf{C})\) 上的向量空间同构。由此推出 D 是弗雷德霍姆自同态，当且仅当空间 \(\mathrm{H}^1 (\mathrm{U};\mathbf{C})\) 是有限维的。然后计算 D 的指标。
\end{enumerate}

